% =====================================================================
% Methodology - material sources (v7.98):
%   * framework/axiom_registry.json v2.0 (taxonomy, 448 entries)
%   * framework/axiom_reduction_sweep1..6_report.md (redemption record)
%   * framework/proof_status.md sec 5 (honest-rate definition),
%     sec 2 (claim tiers + falsifiability criteria), sec 3 (registry)
%   * scripts/compute_proof_rate.py (count_sorry_no_comments = normative
%     reference implementation), scripts/verify_honest_repo.py (CI)
%   * framework/sorry_contamination_audit.md sec 6 (caliber standards)
% \pnum{...} = declared calibers; C6-related figures are canonical
% per p0_reconciliation_report.md (landed 2026-08-21); others are
% verified v7.98 baselines.
% =====================================================================

\subsection{Axiom taxonomy and registry}
\label{sec:taxonomy}
Every \texttt{axiom} declaration in the non-batch corpus is registered
in a machine-readable JSON registry (\texttt{axiom\_registry.json},
schema v2.0) with one of four classes:

\begin{center}
\begin{tabular}{@{}lll@{}}
\toprule
\textbf{Class} & \textbf{Semantics} & \textbf{Disposition} \\
\midrule
\texttt{primitive} & genuinely assumed primitive & keep, declare \\
\texttt{definitional} & restatement of a definition & discharge by
\texttt{rfl}/\texttt{norm\_num}/\texttt{simp} \\
\texttt{schema} & instance of a general transfer pattern &
conditionalize ($P\!\to\!P$) \\
\texttt{placeholder} & deferred proof parked as axiom & full
replacement or explicit debt \\
\bottomrule
\end{tabular}
\end{center}

At the v7.98 baseline the registry contains \pnum{448} entries over
\pnum{116} non-batch files; the \texttt{placeholder} class alone holds
\pnum{42} entries.
% baselines: 448 / 116 / 42 (axiom_registry.json v2.0; D1 review C6,
% sec 1.4-7). NOTE: class-wise full decomposition table pending P0
% (C1 reconciliation: README 478 vs registry 448 vs paper-01 "~350+").
The registry is the single point of truth for redemption accounting
(Sec.~\ref{sec:redemption}); its per-entry provenance (file, line,
sweep that closed it, closing transformation) is what makes the
semantic grading of Sec.~\ref{sec:redemption} auditable rather than
assertive.

\subsection{Redemption methodology and semantic grading}
\label{sec:redemption}
Axioms are discharged in recorded \emph{sweeps} (six at the
reconciled working tree: sweep1--sweep6), each producing a per-entry
report. The methodological
commitment of this paper is that a sweep's output must be graded by
\emph{semantic content added}, not by closure count. Four grades:

\begin{center}
\begin{tabular}{@{}llll@{}}
\toprule
\textbf{Grade} & \textbf{Content added} &
\textbf{Typical transformation} & \textbf{Count} \\
\midrule
G1 vacuous & none & \texttt{axiom T : True} $\to$
\texttt{theorem T : True := trivial} & \pnum{50} \\
G2 trust transfer & none (hypothesis) & \texttt{axiom P} $\to$
\texttt{theorem (h : P) : P := h} & \pnum{90} \\
G3 definitional & definitional equality & data axiom $\to$
\texttt{noncomputable def}; \texttt{rfl} rewrites & \pnum{41} \\
G4 substantive & real proof content & arithmetic / existence proofs
& \pnum{13} \\
\midrule
\multicolumn{3}{@{}l}{\emph{excluded from redemption:} axiom bundling
(kept, declared)} & \pnum{2} \\
\bottomrule
\end{tabular}
\end{center}

% canonical (p0_reconciliation_report.md sec 4, five-class taxonomy
% over sweeps 1-6 = 196 closures; merged 2026-08-21, superseding the
% D1-review C6 estimate and resolving C5/C6):
%   G1 vacuous 50 (sweep1 6 + sweep2 44);
%   G2 trust transfer 90 (sweep1 1 + sweep4 27 + sweep5 30 + sweep6 32);
%   G3 definitional 41 (sweep1 data-def 17 + sweep1 rfl 3 + sweep3 rfl 21);
%   G4 substantive 13 (sweep1 2 + sweep3 10 + sweep4 1);
%   bundling 2 (sweep1 1 + sweep4 1), excluded from redemption.
% C5 resolved: sweep3 ledger is 31 (method table 37 was a miscount).

Two design points deserve emphasis. First, \emph{conditionalization}
(G2) is honest but content-free: it converts an axiom into a theorem
whose hypothesis is the axiom's statement, preserving downstream
elaboration while making the debt visible at every use site. Second,
the grade distribution itself is the deliverable: only \pnum{13}
of \pnum{196} closures are G4. Reporting only ``\pnum{196} axioms
redeemed'' would be the exact failure mode (F2) this methodology
exists to prevent.

% TODO(P1, metatheorem): formalize that G1-G3 transformations yield
% conservative extensions of the ambient theory (D1 review sec 1.4-6).
% Claimed at governance level; not yet a formal theorem.

\subsection{The honest proof-rate metric and caliber discipline}
\label{sec:metric}
The repository's normative definition
(\texttt{proof\_status.md} sec 5):
\begin{quote}
\emph{true proof rate} $=$ (theorems derived from axioms and known
theorems with zero \texttt{postulate}) $/$ (total claimed theorems).
\end{quote}
Operationally, two calibers coexist and both are legitimate when
declared:
\begin{itemize}
\item \caliber{code}: \texttt{sorry} is counted after stripping
comments; files declaring \texttt{axiom}/\texttt{postulate} contribute
zero true proofs. Reference implementation:
\texttt{count\_sorry\_no\_comments} in \texttt{compute\_proof\_rate.py}
--- the single normative implementation, not to be reimplemented ad
hoc.
\item \caliber{independent}: additionally excludes trivially-patterned
proofs and axiom-dependent theorems, yielding the strict ``independent
derivation'' rate.
\end{itemize}
The three-tier caliber standard adopted as policy after the incident
of Sec.~\ref{sec:sorry} is:
\begin{center}
\begin{tabular}{@{}lll@{}}
\toprule
\textbf{Tier} & \textbf{Meaning} & \textbf{Admissible for} \\
\midrule
\caliber{string} & raw substring match, comments included &
nothing (quarantine for humans) \\
\caliber{code} & comment-stripped, code-level pattern &
contamination judgment, CI \\
\caliber{CI-pattern} & anchored CI grep with comment exclusion &
gates, dashboards \\
\bottomrule
\end{tabular}
\end{center}
The rule: \emph{a metric is not a number; it is a (number, caliber)
pair}, and any headline number without a declared caliber is treated
as a defect (cf.\ C3, C4 in Sec.~\ref{sec:conflicts}).

\subsection{The external claim layer}
\label{sec:claims}
Internal formal status and external epistemic status are separated by
a four-tier label, assigned per claim and registered centrally
(\texttt{proof\_status.md}):

\begin{center}
\begin{tabular}{@{}lp{9.2cm}@{}}
\toprule
\textbf{Tier} & \textbf{Standard} \\
\midrule
THEOREM & formal proof machine-checked in-repo, core free of
\texttt{postulate} \\
THEOREM* & mathematical body proven, relying on known theorems not
re-verified in-framework (reference integrity confirmed) \\
CLAIM & reasoning / numerics / citation support, no independent
formalization or review \\
CONJECTURE & speculative, logically coherent, falsifiable \\
\bottomrule
\end{tabular}
\end{center}

Every CLAIM and CONJECTURE carries a \emph{falsifiability criterion}
(e.g., for the fine-structure coincidence $\alpha^{-1}\!\approx\!
n_{CS}\!=\!137$: a first-principles computation of $\alpha^{-1}$ to
$\pm 0.01$ independent of the loop-integer self-tuning parameters).
The registered entries at v7.98 span FLT and Langlands (THEOREM*, via
mathlib), CHSH and the Dedekind unique-factorization delegation
(THEOREM), the stratified Chern-number identity (THEOREM, with
\texttt{\#print axioms} reduced to the standard trio
\texttt{propext}, \texttt{Classical.choice}, \texttt{Quot.sound}),
and the Standard-Model gauge-group, Chern--Simons level, and spectral
action claims (CLAIM, with explicit gap chains).
% source: proof_status.md sec 2 (criteria), sec 3 (registry, v1.1)

\subsection{CI discipline and the reproducibility envelope}
\label{sec:ci}
The honesty layer of the repository is machine-checked
(\texttt{verify\_honest\_repo.py --ci}: forbidden-term scan with
deletion-context allowance, required honesty markers, deleted-claims
ledger with $>$\,10 entries, BibTeX integrity, new-paper presence).
Its \emph{boundary is stated explicitly}: the CI checks the honesty
layer, not Lean proof content --- conflation of the two is itself a
caliber error. The Lean side is pinned: mathlib4 at commit
\texttt{8a178386}, toolchain Lean 4.29.0, with independent compile
logs recording \texttt{\#print axioms} output. The batch generator is
reproducible from its manifest (\pnum{119831} entries, SHA-256
sampling \pnum{500}, zero mismatches at v7.98).
% TODO(P1-8): containerized reproducibility package (compile logs
% currently carry local paths; no CI artifacts yet - D1 review 1.2).
